\section{Software Architecture and Implementation}

The software system is designed as a platform for real-time control and experimental monitoring, 
consisting of a high-precision firmware engine for the ESP32 and a responsive web interface. 
The integration leverages the Bluetooth Low Energy (BLE) protocol to establish the link between 
the embedded hardware and the user's browser.

\section{Firmware Design and Hardware-Level Optimization}

The ESP32 firmware is developed using the ESP-IDF framework to fine control the hardware's dual-core 
capabilities and real-time features. To ensure deterministic performance, specific hardware-level 
configurations are defined in the project using the \textbf{idf menuconfig}.

\begin{itemize}
	\item \textbf{CPU Frequency:} The system is locked at \textbf{240 MHz} to provide the maximum
	      timing resolution for the signal generation loops.

	\item \textbf{Dual-Core Execution:} The ESP32's dual-core capability is used to separate
	      time-critical signal generation from non-deterministic communication.

	\item \textbf{Watchdog Management:} To prevent system resets during continuous signal
	      generation, the Task Watchdog Timer (TWDT) is explicitly disabled for \textbf{Core 1}
	      (\verb|CONFIG_ESP_TASK_WDT_CHECK_IDLE_TASK_CPU1|). This allows the signal task to occupy the
	      CPU cycles without yielding.
\end{itemize}

\subsection{The Real-Time Signal Engine}

The implementation of the signal engine in \verb|signal_controller.cpp| prioritizes speed and timing 
consistency. Rather than relying on standard high-level APIs, the engine utilizes atomic register 
manipulation to perform simultaneous GPIO updates. Key highlights of this implementation include:

\begin{itemize}
\item \textbf{Atomic Switching for Precision:} To protect the voltage converter from 
"shoot-through" currents, the firmware enforces a precise dead-time interval during switching 
transitions. As shown in Listing \ref{lst:switching}, the system utilizes atomic registers 
(\verb|GPIO.out_w1tc| and \verb|GPIO.out_w1ts|) to pull pins low immediately before applying a 
sub-microsecond delay calibrated to the 240 MHz clock. 

\begin{lstlisting}[caption={Atomic Switching and Dead-Time Implementation}, label={lst:switching}] 
// From signal_controller.cpp 
if (clear_mask) { 
    GPIO.out_w1tc = clear_mask; // Atomic clear (Dead Time start)

    uint32_t delay_val = is_rising ? g_cycle_us_delay_up : g_cycle_us_delay_down;
    helper_delay_cycles(delay_val); 

    GPIO.out_w1ts = set_mask;   // Atomic set (Dead Time end)
} 
\end{lstlisting}

	\item \textbf{Double-Buffering Logic:} To allow dynamic parameter updates without halting the 
		hardware output, the firmware maintains an "active" and "inactive" dataset. While 
		\textbf{Core 1} drives the converter using the active set, \textbf{Core 0} parses new 
		commands into the inactive set.

	\item \textbf{Seamless Synchnization:} As illustrated in Listing \ref{lst:buffer}, a synchronization flag 
		(\verb|g_ds_update_pending|) triggers a pointer swap at the start of the next signal cycle, 
		ensuring a glitch-free transition between experimental patterns. 
		

\begin{lstlisting}[caption={Double-Buffer Synchronization Mechanism}, label={lst:buffer}] 
// From signal_controller.cpp 
if (g_ds_update_pending) { 
	if (g_active_set == SignalSet::SET_A) {
		dataset = &g_dataset_b; 
		g_active_set = SignalSet::SET_B; 
	} else { 
		dataset = &g_dataset_a; 
		g_active_set = SignalSet::SET_A; 
	} 

	g_ds_update_pending = false; 
} 
\end{lstlisting}
\end{itemize}

\subsection{Data Acquisition and Calibration}

Real-time monitoring is achieved through a dedicated task that samples the voltages across the 
converter's flying capacitors and the load. Given the non-linear response of the ESP32's 
analog-to-digital converter (ADC), the system incorporates a calibration layer using eFuse-stored 
factory constants and a custom \textbf{Lookup Table (LUT)}.

\begin{lstlisting}[caption={Linear Interpolation Calibration Function}, label={lst:calib}] 
// From helper_analog.cpp 

float esp32_calibration(float value) {
    for (int i = 0; i < calib_numel - 1; ++i) {
        if (value >= calib_from[i] && value <= calib_from[i + 1]) {
            float x1 = calib_from[i];
            float y1 = calib_to[i];
            float x2 = calib_from[i + 1];
            float y2 = calib_to[i + 1];
            return y1 + (value - x1) * (y2 - y1) / (x2 - x1);
        }
    }
    return calib_to[calib_numel - 1];
}
\end{lstlisting}

This calibration utilizes linear interpolation between measured points (Listing \ref{lst:calib}) 
to map raw ADC counts to accurate physical voltages. These values are then packaged into BLE 
notifications for real-time visualization on the web interface.

\subsubsection{Configuration and Control Logic}

The following tables (\ref{tab:system-config} and \ref{tab:ble-commands}) summarize the critical settings and commands required to reproduce the 
system's performance and manage experimental parameters.

\begin{table}[htbp]
	\centering
	\caption{Hardware-Level Firmware Build Configurations}
	\label{tab:system-config}
	\begin{tabular}{llp{8cm}}
		\toprule
		\textbf{Configuration Parameter} & \textbf{Value} & \textbf{Function}                                                        \\
		\addlinespace
		CPU Frequency                    & 240 MHz        & Maximizes timing resolution for helper\_delay\_cycles and system delays. \\
		\addlinespace
		FreeRTOS Cores                   & 2 Cores        & Enables task pinning for core isolation between Signal and BLE domains.  \\
		\addlinespace
		Task Watchdog (Core 1)           & Disabled       & Prevents system reset while the signal task maintains CPU occupancy.     \\
		\addlinespace
		ADC Calibration                  & LUT            & Necessary to correct native non-linearities for analog data acquisition. \\
		\addlinespace
		BLE Stack                        & Bluedroid      & Provides the GATT server infrastructure for remote control.              \\
		\bottomrule
	\end{tabular}
\end{table}


\begin{table}[htbp]
	\centering
	\caption{BLE Communation and Command Router}
	\label{tab:ble-commands}
	\begin{tabular}{llp{8cm}}
		\toprule
		\textbf{Command String} & \textbf{Description} & \textbf{Operational Effect}                                                \\
		\addlinespace
		\verb|SIGNAL:{t;m}|     & Raw Vector Update    & Updates the timing (t) and mode (m) vectors in the inactive dataset.       \\
		\addlinespace
		\verb|SET_ALPHA:x|      & Automatic PSM        & Calculates a Phase-Shifted Modulation pattern for a specific duty cycle x. \\
		\addlinespace
		\verb|START / STOP|     & State Control        & Toggles the signal generation loop between running and idle states.        \\
		\addlinespace
		\verb|STATUS|           & System Metrics       & "Returns current cycle counts, active datasets, and hardware statuses."    \\
		\bottomrule
	\end{tabular}
\end{table}



















\vspace*{20cm}
IK BEN HIER


\subsection{Part 2: The Web Control Interface}

This is the user interface of the system for controlling the ESP32.
It is build with web technologies to run in any browser that supports
Web Bluetooth.

\subsubsection{Core Capabilities}

\begin{itemize}
	\item \textbf{React \& TypeScript:} The application is build using React, a popular
	      JavaScript library for building user interfaces, and TypeScript, which adds static
	      typing to JavaScript. This combination allow a code that is both portable and
	      maintainable.
	\item \textbf{Web Bluetooth API:} The Web Bluetooth API is a browser standard that
	      exposes Bluetooth functionality to JavaScript, allowing the application to
	      \subitem \textbf{Scan} for nearby BLE devices
	      \subitem \textbf{Connect} to a selected device (ESP32)
	      \subitem \textbf{Discover} its services and characteristics
	      \subitem \textbf{Write} data to characteristics (to send commands)
	      \subitem \textbf{Subscribe} to notifications from characteristic (to receive sensor data)
\end{itemize}

This eliminates the need for any native software, drivers, or browser plugins. The entire control
system runs directly in the browser.


\begin{lstlisting}[caption={Bluetooth connect\_device}]
// from bluetooth.ts

export async function connect_device() {
    try {
        let navigatorObject: any = window.navigator;

        // Open the device picker dialog, filtering for our specific service
        g_bluetoothDevice = await navigatorObject.bluetooth.requestDevice({
            filters: [{ services: [SERVICE_UUID] }],
        });

        // Connect to the device's GATT server
        g_gattServer = await g_bluetoothDevice.gatt.connect();

        // Get the service we're interested in
        const service = await g_gattServer.getPrimaryService(SERVICE_UUID);

        // Get the characteristic we'll use for reading/writing data
        g_characteristic = await service.getCharacteristic(CHARACTERISTIC_UUID);

    } catch (error: any) {
        // Handle any errors during the connection process
    }
}
\end{lstlisting}

\subsection{Web Interface}

The web interface is built with React.js to provide an easy user experience.
This interface serves as the primary connection point between the user and the hardware,
linking directly to the ESP32 microcontroller via Bluetooth Low Energy (BLE).

From the web application (Figure \ref{fig:sw_web_interface}), users can easily send string commands to the device, as well as trigger automated
sequences for common experimental procedures. A key function of the interface is its ability to perform
on-the-fly computations to generate the control signals needed for the power switches. It also
receives and processes data from the ESP32's BLE connection, specifically the analog measurements from
the circuit. The received data is then visually represented in a real-time chart, allowing for immediate
analysis of the circuit's performance.


\begin{figure}[h!]
	\centering
	\includegraphics[width=140mm]{images/hw_esp32_pins.pdf}
	\caption{ESP32 Digital and Analog mapping}
	\label{fig:sw_esp_pins}
\end{figure}


\begin{figure}[h!]
	\centering
	\includegraphics[width=140mm]{images/sw_web_interface.pdf}
	\caption{Web Interface to connect and control the ESP32.}
	\label{fig:sw_web_interface}
\end{figure}

The web interface provides single-page application for interacting
with the hardware in real-time. It connects to the ESP32 via Bluetooth Low Energy (BLE) and serves as
a command center for controlling the circuit and analyzing its performance. The following are the key
features of the interface:

\begin{itemize}
	\item \textbf{1. Connect/Disconnect:} Establish or terminate a BLE connection to the ESP32.
	\item \textbf{2. Status command:} Sends the command to get the status of the ESP32 program.
	\item \textbf{3. copy:} Copy the logged data to use in plot or analysis.
	\item \textbf{4. Clear Data:} Clear the logged data.
	\item \textbf{5. $<img>$:} When clicked on the image, the web interface then shows the schematics of the circuit and the connections of the
	      ESP32. This is used as a reference while connecting the circuit.

	\item \textbf{6. Enter command:} Send raw command to ESP32 through BLE.

	\item \textbf{7. SEND:} Send the raw command to the ESP32.

	\item \textbf{8. START:} Commands the ESP32 to start reproducing the signal created.
	\item \textbf{9. STOP (Low):} Commands the ESP32 to set all digital outputs to LOW (0).
	\item \textbf{10. HIGH:} Commands the ESP32 to set all digital outputs to HIGH (1).

	\item \textbf{11. Alpha:} Input tor selecting the alpha value to compute the signal.
	\item \textbf{12. Time:} Time vector in micro seconds ($\mu s$)

	      each element is a amount of $\mu s$ that the mode is going to be executed.

	\item \textbf{13. Mode:} Mode vector for the digital output

	      the mode is the decimal representation of the three digital outputs.
	      The digital output is set to send value to D4, D5 and D6 (figure \ref{fig:sw_esp_pins}).

	      When 13.1 is clicked then the interface shows the decomposed values for each channel D4, D5 and D6.

	      The example in figure \ref{fig:sw_web_interface}, generates the digital signals
	      d4: [0, 0, 0, 1, 1, 1], d5: [0, 1, 1, 1, 0, 0], d6: [1, 1, 0, 0, 0, 1] as result of the
		  signal generated by the Alpha of 0.5.

	      The signal is represented in Figure \ref{fig:sw_signal}.

	      \begin{figure}[h!]
		      \centering
		      \includegraphics[width=140mm]{images/sw_signal.pdf}
		      \caption{Signal for $\alpha$ of 0.4}
		      \label{fig:sw_signal}
	      \end{figure}

	\item \textbf{13.1. UPLOAD:} Shows the Binary representation of the signal modes.

	\item \textbf{14. UPLOAD:} Send the Signal (Time and Mode) to the ESP32.

	\item \textbf{15. Chart Toggle:} Collapses or expand the Chart to show the voltage reading via BLE.
	\item \textbf{16. Advanced Menu Toggle:} Collapses or expand the menu with pre-configured facility commands to send messages via BLE.
	\item \textbf{17. Monitor Chart:} When `Start Listening` is active, then the Monitor chart shows the values of the analog readings from the ESP32.
	\item \textbf{18. System Log:} Displays the log messages from the ESP32.
	\item \textbf{19. Advanced Menu:} Collection of buttons with preconfigured commands.
\end{itemize}

The system's behavior is described by a state diagram in Figure \ref{fig:state_diagram} that shows the
Overall System Flow, which transitions from an initial Setup phase into two main concurrent processes: BLE Task States and Signal Task States.

The BLE (Bluetooth Low Energy) task handles wireless communication
and sensor data, cycling through states such as BLE\_IDLE (when no operations are active), ANALOG\_READING
(when it's actively reading sensor data), and SIGNAL\_READING (when it's processing new signal information).

Concurrently, the Signal task manages the output signals, operating in states like IDLE (with no signals
generated), HIGH\_ALL (where all outputs are forced high), and SIGNAL\_RUN (for active signal pattern
generation). This dual-task structure allows the system to manage sensor inputs and signal outputs independently.

\begin{figure}[h]
	\centering
	\includegraphics[scale=0.5]{images/sw_state_machine.png}
	\caption{State Diagram of the SW Operation Modes}
	\label{fig:state_diagram}
\end{figure}
